\section{The coordinatewise lasso}
\setcounter{equation}{0}

\noindent[k = 4: the statements of this section stood in Sections 3.2 to 3.6 of
the paper, and their proofs in Sections S5 to S7 of this supplement.]

This section gives the theory of the estimator obtained by replacing
$\nGw{\btheta}$ by $\norm{\btheta}_{1}$ in equation (2.5) of the paper, the
lasso of \citet{tibshirani1996regression} with the levels left unpenalized,
which the software offers as an option. Assumptions 1 to 5, Lemmas 1 and 2,
Proposition 2 and Lemmas~\ref{lem:profile}, \ref{lem:columns} and \ref{lem:l1}
are used as they stand, with the notation $\Szero$, $s_0$ and $v$ of
Section~S1. The penalty level is
\begin{equation}
  \label{eq:lasso-lambda}
  \penn \;=\; 2\sigma B_X\sqrt{2C_U}\,
  \sqrt{\frac{2\log(2d_n/\alpha)}{n}} \;\asymp\;
  \sigma B_X \sqrt{C_U}\,\sqrt{\frac{J_n}{n}} ,
\end{equation}
with $d_n$ the number of penalized columns at resolution $J_n$.

\subsection*{Statements}

\begin{theorem}[Oracle inequality for the lasso]
\label{thm:lasso-oracle}
Let Assumption 1 hold, suppose $\Amat$ has rank $p$, and take
$\pen \ge 2\pen_0$ with
$\pen_0 = \sigma\smax\{2\log(2d/\alpha)/n\}^{1/2}$. On the event
$\mathcal{T} = \{\norm{\Btil\T\bveps/n}_{\infty} \le \pen/2\}$, which has
probability at least $1 - \alpha$:
\begin{enumerate}
\item[(i)] with no condition on the design,
$\normn{\Btil v}^{2} \le 6\pen\norm{\thetastar}_{1} + 4\normn{\bb}^{2}$;
\item[(ii)] if $\gtil > 0$,
\[
  \normn{\Btil v}^{2} \le \frac{64\pen^{2}s_0}{\gtil} + 16\normn{\bb}^{2},
  \qquad
  \norm{v}_{2}^{2} \le \frac{64\pen^{2}s_0}{\gtil^{2}}
  + \frac{16\normn{\bb}^{2}}{\gtil} ;
\]
\item[(iii)] in either regime,
$\normn{\widehat{f}-f} \le \normn{\Btil v} + 2\normn{\bb}
+ \normn{\PA\bveps}$, and
$\E(\normn{\PA\bveps}^{2}\mid\bX,\bU) \le \sigma^{2}p/n$.
\end{enumerate}
\end{theorem}

\begin{theorem}[Rate of linear approximation for the lasso]
\label{thm:lasso-linear}
Let Assumptions 1 to 5 hold with $p$, $q$ and the constants fixed, and take
\begin{equation}
  \label{eq:lasso-Jn}
  J_n = \min\bigl\{J \in \mathbb{N} :\;
  2^{J} \ge (n/\log n)^{1/(2\seff+1)}\bigr\},
\end{equation}
so that $J_n = \log_2 n/(2\seff+1) + O(\log\log n)$. Then the requirement of
Proposition 2(ii) holds along $J_n$, being itself of the exact order
$(\log n/n)^{2\seff/(2\seff+1)}$, and, with $\penn$ as in
\eqref{eq:lasso-lambda},
\[
  \normn{\widehat{f}-f}^{2}
  \;=\; O_p\Bigl\{\Bigl(\frac{\log n}{n}\Bigr)^{2\seff/(2\seff+1)}\Bigr\} .
\]
At this resolution the estimation term and the bias term are of the same
order, and the term $\sigma^{2}p/n$ of Theorem~\ref{thm:lasso-oracle}(iii) is
of strictly smaller order.
\end{theorem}

\begin{corollary}[Components and coefficients for the lasso]
\label{cor:lasso-components}
In the setting of Theorem~\ref{thm:lasso-linear}, and with $r_n$ denoting the
rate of that theorem,
\[
  \sum_{\cv,\md}\bigl\lVert\widehat{g}_{\cv\md}-\gcomp{\cv}{\md}
  \bigr\rVert_{L_2[0,1]}^{2} = O_p(r_n),
  \qquad
  \bigl\lVert\widehat{\fcoef{\cv}}-\fcoef{\cv}
  \bigr\rVert_{L_2([0,1]^{q})} = O_p(r_n^{1/2})
\]
for $\cv = 1,\dots,p$, and the same statements hold with the rate of
Theorem~\ref{thm:lasso-main} in place of $r_n$ under the conditions of that
theorem.
\end{corollary}

% E5e (E1.14; D52): a inflacao de penn acima do teto, eq. mu de 05-taxas.tex.
The level $\penn$ is admissible only while $2^{J_n}\log d_n/n \to 0$, which
Lemma~\ref{lem:columns} needs for $\smax^{2} \le 2B_X^{2}C_U$; above that
ceiling $\smax^{2}$ grows with $2^{J}/n$. With $c_{\psi} =
7\norm{\psi}_{\infty}^{2}/(12C_U)$ as in Section~S6, write
\begin{equation}
  \label{eq:lasso-lamplus}
  \mu_{J} = 1 + c_{\psi}\frac{2^{J}\log d}{n},
  \qquad
  \penn^{+} = \mu_{J_n}^{1/2}\penn ;
\end{equation}
when $2^{J_n}\log d_n/n \to 0$, $\penn^{+} = \penn\{1 + o(1)\}$.

% derivations/05-taxas.tex, Proposicao 4, emendada por D49(a), com os itens
% (iii) e (iv) de E1.14 (E5e).
\begin{proposition}[Slow rate for the lasso, no design condition]
\label{prop:lasso-slow}
% E5f (D54): as cotas superiores das Assumptions 2 e 3, que o penn e a prova usam.
Let Assumptions 1, 4 and 5 hold, together with the upper
bounds in Assumptions 2 and 3, with $\Amat$ of rank $p$ and $\penn$ as in
\eqref{eq:lasso-lambda}. Then
$\normn{\widehat{f}-f}^{2} = O_p\{\penn\norm{\thetastar}_{1} +
\mathcal{B}_{J_n}^{2} + p/n\}$, where
$\norm{\thetastar}_{1} \le pqC_g\sum_{\lev<J_n}2^{\lev(1/2-s)}$ does not
depend on $\pi$, whenever $2^{J_n}\log d_n/n \to 0$; with $\penn^{+}$ of
\eqref{eq:lasso-lamplus} in place of $\penn$, in the estimator and in the
bound, the same holds for every sequence $J_n$. In particular, (i) if $s > 1/2$ and $\seff > 1/4$ then
$\normn{\widehat{f}-f}^{2} = O_p\{(J_n/n)^{1/2}\}$ for every
$J_n \ge \log_2 n/(4\seff)$; and (ii) if $s < 1/2$ and $\seff > s/2$, the choice
$2^{J_n} \asymp (n/J_n)^{1/(1-2s+4\seff)}$ gives
$\normn{\widehat{f}-f}^{2} = O_p\{(J_n/n)^{2\seff/(1-2s+4\seff)}\}$, an
exponent that equals $2s/(2s+1)$ when $\pi \ge 2$. With $\penn^{+}$:
\begin{enumerate}
\item[(iii)] (the pairs that (i) and (ii) exclude, all with $\pi < 2$)
(a) if $s > 1/2$ and $\seff \le 1/4$, the choice
$2^{J_n} \asymp (n/\log n)^{2/(1+4\seff)}$ gives
$\normn{\widehat{f}-f}^{2} = O_p\{(\log n/n)^{4\seff/(1+4\seff)}\}$;
(b) if $s < 1/2$ and $\seff \le s/2$, the choice
$2^{J_n} \asymp (n/\log n)^{1/(1-s+2\seff)}$ gives
$\normn{\widehat{f}-f}^{2} = O_p\{(\log n/n)^{2\seff/(1-s+2\seff)}\}$;
in both, $2^{J_n} \gtrsim n/\log n$, above the ceiling of
Lemma~\ref{lem:columns}, and the rates join those of (i) and (ii) at
$\seff = 1/4$ and $\seff = s/2$;
\item[(iv)] (the boundary $s = 1/2$, with $\pi > 1$)
$\norm{\thetastar}_{1} \le pqC_gJ_n$, and: if $\seff \ge 1/4$, that is
$\pi \ge 4/3$, the choice $2^{J_n} \asymp (n/\log n)^{1/(4\seff)}$ gives
$\normn{\widehat{f}-f}^{2} = O_p\{n^{-1/2}(\log n)^{3/2}\}$; if
$\seff < 1/4$, that is $1 < \pi < 4/3$, the choice
$2^{J_n} \asymp \{n/(\log n)^{2}\}^{2/(1+4\seff)}$ gives
$\normn{\widehat{f}-f}^{2} = O_p[\{(\log n)^{2}/n\}^{4\seff/(1+4\seff)}]$.
\end{enumerate}
\end{proposition}

% E5f (E1.15; D54): derivations/05-taxas.tex, Proposicao 8 (numeracao global;
% prop:truncada); a prova vem depois da Remark S8.2.
Proposition~\ref{prop:lasso-slow} reads Theorem~\ref{thm:lasso-oracle}(i) at
$\thetastar$ and pays $\pen|\theta^{*}_{a}|$ also for the coefficients
smaller than $\pen$, which would cost less as bias. Lemma~\ref{lem:comparator}
below holds for every comparator on the same event, and the comparator
truncated at $t > 0$,
$\thbar_{a} = \theta^{*}_{a}\mathbf{1}\{|\theta^{*}_{a}| > t\}$ for
$a = 1,\dots,d$, charges the penalty only for the coefficients above $t$. What
remains is
\begin{equation}
  \label{eq:R1}
  \Rone(\btheta;t) \;=\; \sum_{a=1}^{d}\min\bigl(t|\theta_{a}|,\,\theta_{a}^{2}\bigr)
  \;\le\; t\norm{\btheta}_{1},
\end{equation}
the counterpart in the $\ell_1$ norm of the ideal risk
$\sum_{a}\min(\theta_{a}^{2},t^{2})$.

\begin{proposition}[Slow rate for the lasso through the truncated comparator]
\label{prop:lasso-truncated}
Let the assumptions of Proposition~\ref{prop:lasso-slow} hold.
\begin{enumerate}
\item[(i)] On the event $\mathcal{T}$, for every $\pen \ge 2\pen_0$ and
every $t > 0$, with $\thbar$ truncated at $t$,
\[
  \normn{\widehat{f}-f}^{2}
  \;\le\; 18\pen\norm{\thbar}_{1} + 48\normn{\Bmat(\thetastar-\thbar)}^{2}
  + 48\normn{\bb}^{2} + 3\normn{\PA\bveps}^{2} .
\]
If $t$ is deterministic,
$\E\normn{\Bmat(\thetastar-\thbar)}^{2} \le
\kappa_2C_U\norm{\thetastar-\thbar}_{2}^{2}$. With $t = \pen$ deterministic,
$18\pen\norm{\thbar}_{1} + 48\kappa_2C_U\norm{\thetastar-\thbar}_{2}^{2} \le
\max(18,48\kappa_2C_U)\Rone(\thetastar;\pen)$, and
\[
  \normn{\widehat{f}-f}^{2} \;=\; O_p\bigl\{\Rone(\thetastar;\penn)
  + \mathcal{B}_{J_n}^{2} + p/n\bigr\}
\]
whenever $2^{J_n}\log d_n/n \to 0$; with $\penn^{+}$ of
\eqref{eq:lasso-lamplus} in place of $\penn$, in the estimator and in the
bound, the same holds for every sequence $J_n$ with $2^{J_n} \le n^{2}$.
\item[(ii)] For every $J$ and every $t > 0$, with $\tau = (s+1/2)^{-1}$: if
$s < 1/2$,
\begin{equation}
  \label{eq:count-lasso-truncated}
  \Rone(\thetastar;t) \;\le\; pq\,C_g^{\tau}\,t^{4s/(2s+1)}
  \Bigl\{\frac{1}{1-2^{s-1/2}} + \frac{1}{1-2^{-\min(\pi,2)\seff}}\Bigr\};
\end{equation}
if $s \ge 1/2$, nothing is gained in order: $\Rone(\thetastar;t) \le
t\norm{\thetastar}_{1}$, and for $t \le C_g$ the class of Assumption~5
contains a $\thetastar$ with $\Rone(\thetastar;t) \ge tC_g$ if $s > 1/2$,
and one with $\Rone(\thetastar;t) \ge
tC_g\min\{J,\,\lfloor\log_2(C_g/t)\rfloor + 1\}$ if $s = 1/2$, the order of
$t\norm{\thetastar}_{1}$ in Proposition~\ref{prop:lasso-slow}.
\item[(iii)] If $s < 1/2$ and $\seff > s/(2s+1)$, then with
$J_n = \lceil c\log_2 n\rceil$ and $c$ in the window \eqref{eq:lasso-window},
\[
  \normn{\widehat{f}-f}^{2}
  \;=\; O_p\Bigl\{\Bigl(\frac{\log n}{n}\Bigr)^{2s/(2s+1)}\Bigr\},
\]
the prediction rate of Theorem~\ref{thm:lasso-main}.
\item[(iv)] If $s < 1/2$ and $\seff \le s/(2s+1)$, which forces $\pi < 2$,
then with $\penn^{+}$ and the choice
$2^{J_n} \asymp (n/\log n)^{2s/\{s+\seff(2s+1)\}}$, which has
$2^{J_n} \gtrsim n/\log n$,
\[
  \normn{\widehat{f}-f}^{2}
  \;=\; O_p\Bigl\{\Bigl(\frac{\log n}{n}\Bigr)^{4s\seff/\{s+\seff(2s+1)\}}\Bigr\};
\]
at $\seff = s/(2s+1)$ the exponent is $2s/(2s+1)$, that of (iii).
\end{enumerate}
\end{proposition}

Proposition~\ref{prop:lasso-truncated} is a statement on prediction: it uses
the design as Proposition~\ref{prop:lasso-slow} does, and, in the new bias,
the upper bound $\lmaxx(\Gram) \le \kappa_2C_U$; the components still need
the design condition (Corollary~\ref{cor:lasso-components}). When $s < 1/2$
and $\pi < 2$ it improves on Proposition~\ref{prop:lasso-slow} in every pair:
(iii) gives the rate of Theorem~\ref{thm:lasso-main} in the window of that
theorem, with a gain of $2(s-\seff)(1-2s)/\{(2s+1)(1-2s+4\seff)\}$ in the
exponent over Proposition~\ref{prop:lasso-slow}(ii), and (iv) replaces
Proposition~\ref{prop:lasso-slow}(iii)(b) and the part of
Proposition~\ref{prop:lasso-slow}(ii) with $s/2 < \seff \le s/(2s+1)$. When
$\pi \ge 2$ the rate is that of Proposition~\ref{prop:lasso-slow}(ii), now
for every $c \in [1/(2s+1),1)$. When $s \ge 1/2$ nothing changes, for the
reason given after Proposition~\ref{prop:truncated}.

For a vector $\btheta$ let $\theta_{(1)} \ge \theta_{(2)} \ge \cdots$ be the
nonincreasing rearrangement of the absolute values of its entries and let
$\btheta^{(M)}$ retain the $M$ largest of them.

\begin{lemma}[Besov regularity implies compressibility]
\label{lem:lasso-weak}
Under Assumptions 4 and 5 with $\pi(s+1/2) > 1$, the oracle satisfies, for
every $J$ and every $k \ge 1$,
\[
  \theta^{*}_{(k)} \;\le\; C_\tau\,k^{-1/\tau},
  \qquad \tau = \frac{1}{s+1/2},
  \qquad C_\tau = C_g(pq\,A_{s,\pi})^{s+1/2},
  \qquad A_{s,\pi} = 1 + \frac{1}{1-2^{1-\pi(s+1/2)}},
\]
with constants that do not depend on $J$ and with $\tau \in (0,2)$ for every
$s > 0$. Consequently
$\sum_{k>M}\theta^{*2}_{(k)} \le \tau(2-\tau)^{-1}C_\tau^{2}M^{1-2/\tau}$ for
every $M \ge 1$.
\end{lemma}

The inclusion of Besov balls in weak $\ell_\tau$ balls is classical
\citep[Proposition~9.15]{johnstone2019gaussian}; what the lemma adds is the
constant $C_\tau$, uniform in $J$, for the oracle of the product design, which
is the form in which the proof of Theorem~\ref{thm:lasso-main} uses it.

\begin{theorem}[Rate of nonlinear approximation for the lasso]
\label{thm:lasso-main}
Let Assumptions 1 to 5 hold with $p$, $q$ and all the constants fixed as $n$
grows, and set $\tau = (s+1/2)^{-1}$. Suppose $\seff > s/(2s+1)$ and take
$J_n = \lceil c \log_2 n\rceil$ with
\begin{equation}
  \label{eq:lasso-window}
  \max\Bigl\{\frac{\tau}{2},\; \frac{2-\tau}{4\seff}\Bigr\}
  \;\le\; c \;<\; 1 ,
\end{equation}
a window that is nonempty precisely under that condition. Then, with $\penn$
as in \eqref{eq:lasso-lambda},
\[
  \normn{\widehat{f}-f}^{2}
  \;=\; O_p\Bigl\{\Bigl(\frac{\log n}{n}\Bigr)^{2s/(2s+1)}\Bigr\},
  \qquad
  \sum_{\cv,\md}\bigl\lVert \widehat{g}_{\cv\md}-\gcomp{\cv}{\md}
  \bigr\rVert_{L_2[0,1]}^{2}
  \;=\; O_p\Bigl\{\Bigl(\frac{\log n}{n}\Bigr)^{2s/(2s+1)}\Bigr\} .
\]
\end{theorem}

The logarithmic factor is the price of the single penalty level
\eqref{eq:lasso-lambda}, which takes a union bound over the $d_n$ columns;
Theorem~1 of the paper removes the factor $(\log n)^{2\seff/(2s+1)}$ of it.
Write $\rho_n = (\log n/n)^{\seff/(2\seff+1)}$ for the square root of the rate
of Theorem~\ref{thm:lasso-linear}.

\begin{corollary}[Block structure by thresholding the lasso]
\label{cor:lasso-threshold}
Let the assumptions of Theorem~\ref{thm:lasso-linear} hold, with $J_n$ and
$\penn$ as there.
\begin{enumerate}
\item[(i)] (Screening.) If the thresholds satisfy $t_n/\rho_n \to \infty$,
then $\Prob\{\widehat{\mathcal{S}}(t_n) \subseteq \mathcal{S}\} \to 1$.
\item[(ii)] (Recovery.) If, in addition, Assumption 6 holds with $\rho_n$ in
place of $\rhoG$ and $\limsup_n t_n/\nu_{\min,n} < 1$, then
$\Prob\{\widehat{\mathcal{S}}(t_n) = \mathcal{S}\} \to 1$. Such thresholds
always exist; one is $t_n = (\rho_n\nu_{\min,n})^{1/2}$.
\item[(iii)] (A window at fixed $n$.) With $\alpha$ the level in
\eqref{eq:lasso-lambda}, for every $n$ and every $\alpha' \in (0,1)$ there is
an event, of probability at least $1 - \alpha - \alpha' - \omega_n$ with
$\omega_n \to 0$ not depending on the components, on which
$\widehat{\mathcal{S}}(t) = \mathcal{S}$ for every $t$ with
$\bar{\Delta}_n \le t < \nu_{\min,n} - \bar{\Delta}_n$; here
$\bar{\Delta}_n \le c(\alpha,\alpha')\rho_n$ is the explicit bound
\eqref{eq:Dn} on $\max_{(\cv,\md)}\lVert \widehat{g}_{\cv\md} -
\gcomp{\cv}{\md}\rVert_{L_2[0,1]}$.
\end{enumerate}
\end{corollary}

Parts (i) and (ii) hold also at the resolution of
Theorem~\ref{thm:lasso-main}, with $\rho_n$ replaced by
$\tilde\rho_n = (\log n/n)^{s/(2s+1)}$, the square root of the rate of that
theorem, both in the condition on $t_n$ and in Assumption~6. Since
$\tilde\rho_n \le \rho_n$, with a gap in the exponent when $\pi < 2$, the
separation required is weaker exactly where the components are spatially
inhomogeneous.

\subsection*{Proofs}

\begin{lemma}[Calibration of the penalty]
\label{lem:calib}
Under Assumption 1, for every $\alpha \in (0,1)$,
$\Prob(\norm{\Btil\T\bveps/n}_{\infty} > \pen_0) \le \alpha$ with
$\pen_0 = \sigma\smax\{2\log(2d/\alpha)/n\}^{1/2}$.
\end{lemma}

\begin{proof}
Condition on the covariates, so that $\Btil$ and $\smax$ are fixed and the
errors are independent and sub-Gaussian with parameter $\sigma$. For a column
$a$, $(\Btil\T\bveps)_a/n$ is sub-Gaussian with parameter
$\sigma\lVert\Btil^{(a)}\rVert_2/n$, and, $\MA$ being an orthogonal
projection, $\lVert\Btil^{(a)}\rVert_2 \le \lVert\Bmat^{(a)}\rVert_2 =
\sqrt{n}\,(\Gramhat)_{aa}^{1/2} \le \sqrt{n}\,\smax$. With
$t = \pen_0$ the sub-Gaussian tail bound is at most $\alpha/d$ per column; a
union bound over the $d$ columns and the tower property finish the proof.
\end{proof}

\subsection*{Proof of Theorem~\ref{thm:lasso-oracle}}

\emph{Step 1, the basic inequality.} By Lemma S5.1, $\thetahat$ minimizes the
residualized objective; comparing its value at $\thetahat$ with its value at
$\thetastar$, and using $\MA\Amat = 0$, so that
$\MA\bY = \Btil\thetastar + \MA(\bb+\bveps)$, gives, with
$r = \MA(\bb+\bveps)$,
\begin{equation}
  \label{eq:basic}
  \normn{\Btil v}^2 + 2\pen\norm{\thetahat}_1
  \;\le\; 2\inner{r}{\Btil v}_n + 2\pen\norm{\thetastar}_1 .
\end{equation}

\emph{Step 2, the two noises.} Since $\MA\Btil = \Btil$ and $\MA$ is
symmetric, $\inner{r}{\Btil v}_n = \inner{\bb+\bveps}{\Btil v}_n$. The
stochastic part is handled by duality: on $\mathcal{T}$,
$2|\inner{\bveps}{\Btil v}_n| \le
2\norm{\Btil\T\bveps/n}_{\infty}\norm{v}_1 \le \pen\norm{v}_1$. The
deterministic part is handled by the Cauchy-Schwarz and Young inequalities,
and not by duality, because $\bb$ is not small coordinatewise:
$2|\inner{\bb}{\Btil v}_n| \le 2\normn{\bb}\normn{\Btil v} \le
\normn{\Btil v}^2/2 + 2\normn{\bb}^2$. Substituting both in
\eqref{eq:basic},
\begin{equation}
  \label{eq:basic2}
  \tfrac12\normn{\Btil v}^2 + 2\pen\norm{\thetahat}_1
  \;\le\; \pen\norm{v}_1 + 2\pen\norm{\thetastar}_1 + 2\normn{\bb}^2 .
\end{equation}

\emph{Step 3, the slow rate.} Bounding
$\norm{v}_1 \le \norm{\thetahat}_1+\norm{\thetastar}_1$ in
\eqref{eq:basic2} gives
$\tfrac12\normn{\Btil v}^2 + \pen\norm{\thetahat}_1 \le
3\pen\norm{\thetastar}_1 + 2\normn{\bb}^2$, which is part (i). No property of
$\Gramhat$ was used.

\emph{Step 4, the cone.} Since $\thetastar$ vanishes off $\Szero$,
$\norm{\thetahat}_1 \ge \norm{\thetastar}_1 - \norm{v_{\Szero}}_1 +
\norm{v_{\Szero^c}}_1$ and
$\norm{v}_1 = \norm{v_{\Szero}}_1+\norm{v_{\Szero^c}}_1$; substituting in
\eqref{eq:basic2} and cancelling $2\pen\norm{\thetastar}_1$,
\begin{equation}
  \label{eq:cone}
  \tfrac12\normn{\Btil v}^2 + \pen\norm{v_{\Szero^c}}_1
  \;\le\; 3\pen\norm{v_{\Szero}}_1 + 2\normn{\bb}^2 .
\end{equation}

\emph{Step 5, two cases.} If $2\normn{\bb}^2 \le \pen\norm{v_{\Szero}}_1$,
then \eqref{eq:cone} gives
$\tfrac12\normn{\Btil v}^2 + \pen\norm{v_{\Szero^c}}_1 \le
4\pen\norm{v_{\Szero}}_1$; by the Cauchy-Schwarz inequality and
\eqref{eq:schur}, $\norm{v_{\Szero}}_1 \le \sqrt{s_0}\norm{v}_2 \le
\sqrt{s_0}\normn{\Btil v}/\sqrt{\gtil}$, so
$\normn{\Btil v} \le 8\pen\sqrt{s_0/\gtil}$ and
$\norm{v}_2^2 \le \normn{\Btil v}^2/\gtil \le 64\pen^2s_0/\gtil^2$; the same
display gives $\norm{v_{\Szero^c}}_1 \le 4\norm{v_{\Szero}}_1$ and hence
$\norm{v}_1 \le 5\norm{v_{\Szero}}_1 \le 40\pen s_0/\gtil$. If instead
$2\normn{\bb}^2 > \pen\norm{v_{\Szero}}_1$, then \eqref{eq:cone} gives
$\tfrac12\normn{\Btil v}^2 < 8\normn{\bb}^2$, that is
$\normn{\Btil v}^2 < 16\normn{\bb}^2$, together with
$\pen\norm{v}_1 < 10\normn{\bb}^2$ and
$\norm{v}_2^2 < 16\normn{\bb}^2/\gtil$. Each quantity is bounded by the
maximum of the two cases, hence by their sum, which is part (ii).

\emph{Step 6, prediction.} By \eqref{eq:decomp},
$\widehat{f}-f = \Btil v + \PA(\bb+\bveps) - \bb$, and $\PA$ being a
projection, $\normn{\PA\bb} \le \normn{\bb}$, which gives the triangle
inequality of part (iii). Conditionally on the covariates the errors are
independent, centered and of variance at most $\sigma^2$, so
\[
  \E(\normn{\PA\bveps}^2\mid\bX,\bU)
  = \frac1n\tr\{\PA\,\mathrm{Cov}(\bveps\mid\bX,\bU)\}
  \le \frac{\sigma^2}{n}\tr(\PA) \le \frac{\sigma^2p}{n},
\]
since the diagonal entries of $\PA$ are nonnegative and its trace is the rank
of $\Amat$. \qed

\begin{remark}[Why there is no cone condition]
\label{rem:nocone}
The usual treatment of unpenalized coordinates \citep[Section
6.9]{buhlmann2011statistics} places them inside the support and requires
compatibility on the cone $\norm{\delta_{S^c}}_1 \le 3\norm{\delta_S}_1$.
Profiling removes them from the problem, and \eqref{eq:schur} transfers the
full smallest eigenvalue of $\Gramhat$ to the residualized block, so that
Proposition 2(ii) is used as it stands and the constant does not depend on
$\Szero$. If only the restricted version of the design condition were
available, the cone would return, and with the constant $4$ rather than $3$,
since Step 5 produces $\norm{v_{\Szero^c}}_1 \le 4\norm{v_{\Szero}}_1$: the
bias consumes part of the slack.
\end{remark}

\begin{corollary}[The two rates with the design condition instantiated]
\label{cor:instant}
Add Assumptions 2 to 5 and suppose $2^J\log d/n \to 0$, which is implied by
the requirement of Proposition 2(ii). Take
$\pen = 2\sigma B_X\sqrt{2C_U}\{2\log(2d/\alpha)/n\}^{1/2}$. Then, with
probability at least $1-\alpha-\alpha'-\alpha''$ up to terms that tend to
zero,
\[
  \normn{\widehat{f}-f}^2 \;\lesssim\;
  \frac{\sigma^2B_X^2C_U}{\kappa_1c_U}\,\frac{s_0\log(pq\NJ)}{n}
  \;+\; \frac{\mathcal{B}^2_J}{\alpha'} \;+\; \frac{\sigma^2p}{n\alpha''},
\]
and, with no condition on the design, the same bound with
$\pen\norm{\thetastar}_1$ of Lemma S5.4 in place of the first term.
\end{corollary}

\begin{proof}
Intersect four events: $\mathcal{T}$, of probability at least $1-\alpha$, by
Lemma~\ref{lem:calib}; $\{\smax^2 \le 2B_X^2C_U\}$, by Lemma S5.3, which makes the stated
$\pen$ admissible, that is $\pen \ge 2\pen_0$;
$\{\lmin(\Gramhat)\ge\gamma/2\}$, by Proposition 2(ii), which with
\eqref{eq:schur} gives $\gtil \ge \gamma/2$; and
$\{\normn{\bb}^2 \le \mathcal{B}_J^2/\alpha'\}$, which is Markov's inequality
applied to $\E\normn{\bb}^2 = \norm{f-f_J}^2_{L_2(P)} \le \mathcal{B}_J^2$,
the identity holding because the observations are identically distributed.
Substituting in parts (ii) and (iii) of Theorem~\ref{thm:lasso-oracle} gives the statement, and
part (i) with Lemma S5.4 gives the unconditional version.
\end{proof}

The first lemma is Theorem~\ref{thm:lasso-oracle} read at an arbitrary comparator. It consumes no
new step: the proof of that theorem uses the oracle only through the
decomposition of $\bY$ and through the size of its support.

\begin{lemma}[Theorem~\ref{thm:lasso-oracle} with an arbitrary comparator]
\label{lem:comparator}
Under Assumption 1, with $\Amat$ of rank $p$ and $\pen \ge 2\pen_0$, on the
event $\mathcal{T}$ the following hold simultaneously for every
$\bar\bbeta = (\bar\bc,\thbar)$, with $\bbar = f - \bZ\bar\bbeta$,
$\bar M = |\supp(\thbar)|$ and $\bar v = \thetahat-\thbar$:
$\tfrac12\normn{\Btil\bar v}^2 + \pen\norm{\thetahat}_1 \le
3\pen\norm{\thbar}_1 + 2\normn{\bbar}^2$; if $\gtil>0$,
$\normn{\Btil\bar v}^2 \le 64\pen^2\bar M/\gtil + 16\normn{\bbar}^2$ and
$\norm{\bar v}_2^2 \le 64\pen^2\bar M/\gtil^2 +
16\normn{\bbar}^2/\gtil$; and
\begin{equation}
  \label{eq:oracle-min}
  \normn{\widehat{f}-f}^2 \;\le\;
  3\min_{\bar\bbeta}\Bigl\{\frac{64\pen^2\bar M}{\gtil}
  + 20\normn{\bbar}^2\Bigr\} + 3\normn{\PA\bveps}^2 .
\end{equation}
\end{lemma}

\begin{proof}
Follow the proof of Theorem~\ref{thm:lasso-oracle} with
$(\bar\bc,\thbar,\bbar,\supp(\thbar),\bar M,\bar v)$ in place of
$(\bc,\thetastar,\bb,\Szero,s_0,v)$. Step 1 uses only
$\bY = \Amat\bar\bc+\Bmat\thbar+\bbar+\bveps$, which is the definition of
$\bbar$, and $\MA\Amat = 0$. Step 2 uses $\mathcal{T}$, which does not depend
on the comparator, and $\normn{\MA\bbar} \le \normn{\bbar}$. Steps 3 to 5 are
algebra with $\supp(\thbar)$ in place of $\Szero$, using
$\thbar_{\supp(\thbar)^c} = 0$, the Cauchy-Schwarz inequality and
\eqref{eq:schur}. Step 6 gives
$\widehat{f}-\bZ\bar\bbeta = \Btil\bar v + \PA(\bbar+\bveps)$ and hence
$\widehat{f}-f = \Btil\bar v+\PA(\bbar+\bveps)-\bbar$. Since $\mathcal{T}$,
$\gtil$ and the condition on $\pen$ do not involve $\bar\bbeta$, the
inequalities hold simultaneously on the same event and the minimum may be
taken; \eqref{eq:oracle-min} follows from
$(x+y+z)^2 \le 3(x^2+y^2+z^2)$.
\end{proof}

\subsection*{Proof of Lemma~\ref{lem:lasso-weak}}

\emph{The envelope.} Fix $\epsilon>0$ and count
$N(\epsilon) = \#\{a : |\theta^{*}_a| > \epsilon\}$. Inside a block and a
level, Markov's inequality in the $\ell_\pi$ norm gives
$\#\{\tsl : |\coefs{\cv}{\md}{\lev}{\tsl}|>\epsilon\} \le
\epsilon^{-\pi}\norm{\theta^{*}_{\cv\md,\lev\cdot}}^{\pi}_{\pi}$, and
Assumption 5 gives
$\norm{\theta^{*}_{\cv\md,\lev\cdot}}^{\pi}_{\pi} \le
C_g^{\pi}2^{-\lev\pi(s+1/2-1/\pi)} = C_g^{\pi}2^{\lev(1-a)}$ with
$a = \pi(s+1/2)>1$. Since a level has $2^{\lev}$ coordinates,
\[
  N_{\cv\md}(\epsilon) \;\le\; \sum_{\lev\ge0}
  \min\{2^{\lev},\, C_g^{\pi}\epsilon^{-\pi}2^{\lev(1-a)}\} .
\]
The two arguments agree at $2^{\lev_\epsilon a} = (C_g/\epsilon)^{\pi}$, that
is at $2^{\lev_\epsilon} = (C_g/\epsilon)^{\tau}$, since
$\pi/a = (s+1/2)^{-1} = \tau$. Below that level the minimum is $2^{\lev}$ and
the sum is at most $2^{\lev_\epsilon}$; above it the minimum is the second
term, which decreases geometrically with ratio $2^{1-a}<1$, and the sum is
$2^{\lev_\epsilon}/(1-2^{1-a})$. Adding the two parts and the $pq$ blocks
gives $N(\epsilon) \le pqA_{s,\pi}(C_g/\epsilon)^{\tau}$; truncating at
$\lev<J$ only decreases the count, so the bound holds for every $J$. Finally,
$\epsilon < \theta^{*}_{(k)}$ implies $N(\epsilon) \ge k$, hence
$\epsilon \le C_g(pqA_{s,\pi}/k)^{1/\tau}$, and letting $\epsilon$ increase to
$\theta^{*}_{(k)}$ gives the envelope; and $\tau = (s+1/2)^{-1} < 2$ for every
$s>0$.

\emph{The tail.} With $2/\tau>1$,
$\sum_{k>M}\theta^2_{(k)} \le C_\tau^2\sum_{k>M}k^{-2/\tau} \le
C_\tau^2\int_M^{\infty}x^{-2/\tau}dx = C_\tau^2M^{1-2/\tau}\tau/(2-\tau)$.
\qed

\subsection*{Proof of Proposition~\ref{prop:lasso-slow}}

By parts (i) and (iii) of Theorem~\ref{thm:lasso-oracle}, on $\mathcal{T}$,
$\normn{\Btil v}^2 \le 6\pen\norm{\thetastar}_1+4\normn{\bb}^2$ and
$\normn{\widehat{f}-f}^2 \le 3(\normn{\Btil v}^2+4\normn{\bb}^2+
\normn{\PA\bveps}^2)$, and no property of $\Gramhat$ enters. Lemma S5.3 makes
$\penn$ admissible with probability tending to one, since
$2^{J_n}\log d_n/n \to 0$ in both regimes, and it is for this that the conditions
$\seff > 1/4$ in (i) and $\seff > s/2$ in (ii) are needed: without them the choice
in (i) gives $2^{J_n} \ge n^{1/(4\seff)} \ge n$, and the balance in (ii) gives
$2^{J_n} \gtrsim n/\log n$, and parts (iii) and (iv) cover those pairs with
$\penn^{+}$; Markov's inequality
applied to $\E\normn{\bb}^2 \le \mathcal{B}^2_{J_n}$, from Lemma 2, and to
$\E(\normn{\PA\bveps}^2\mid\bX,\bU) \le \sigma^2p/n$ gives the general
statement. If $s>1/2$, Lemma S5.4 gives $\norm{\thetastar}_1 = O(1)$ and
$\penn\norm{\thetastar}_1 \asymp (J_n/n)^{1/2}$, and it suffices that the bias
not dominate, that is $2^{-2J_n\seff} \le (J_n/n)^{1/2}$, which holds for
$J_n \ge \log_2(n/J_n)/(4\seff)$. If $s<1/2$, then
$\norm{\thetastar}_1 \asymp 2^{J_n(1/2-s)}$, the two terms are
$(J_n/n)^{1/2}2^{J_n(1/2-s)}$ and $2^{-2J_n\seff}$, and equating them gives
$2^{J_n} \asymp (n/J_n)^{1/(1-2s+4\seff)}$ with common value
$(J_n/n)^{2\seff/(1-2s+4\seff)}$; when $\seff = s$ the denominator is
$1+2s$.

% E5e (E1.14; D52): a prova dos itens (iii) e (iv) da Proposicao 4 de 05-taxas.tex.
\emph{The level above the ceiling, for (iii) and (iv).} By
\eqref{eq:smax-ceiling}, $\Prob(\smax^{2} > 2B_X^{2}C_U\mu_{J}) \le 1/d$ for
every $J$. Off that event,
\[
  2\pen_0 \;\le\; 2\sigma(2B_X^{2}C_U\mu_{J})^{1/2}
  \Bigl\{\frac{2\log(2d/\alpha)}{n}\Bigr\}^{1/2} \;=\; \penn^{+},
\]
and the argument of the general statement goes through with
$\penn^{+}$ and no condition on $J_n$: on $\mathcal{T}$,
$\normn{\widehat{f}-f}^{2} \le 18\penn^{+}\norm{\thetastar}_{1} +
24\normn{\bb}^{2} + 3\normn{\PA\bveps}^{2}$, and Markov's inequality
concludes. In every choice of (iii) and (iv), $2^{J_n} \le n^{2}$, so that
$\log d_n \asymp J_n \asymp \log n$ and
\begin{equation}
  \label{eq:lam-plus}
  \penn^{+} \;\asymp\; \Bigl(\frac{\log n}{n}\Bigr)^{1/2}
  \Bigl(1 + \frac{2^{J_n}\log n}{n}\Bigr)^{1/2},
\end{equation}
in particular $\penn^{+} \asymp 2^{J_n/2}\log n/n$ if
$2^{J_n} \gtrsim n/\log n$.

\emph{Part (iii)(a).} With $s > 1/2$, $\norm{\thetastar}_{1} = O(1)$ as in
(i). With $2^{J_n} \asymp (n/\log n)^{2/(1+4\seff)} \gtrsim n/\log n$,
because $\seff \le 1/4$, \eqref{eq:lam-plus} gives
$\penn^{+}\norm{\thetastar}_{1} \asymp (\log n/n)(n/\log n)^{1/(1+4\seff)} =
(\log n/n)^{4\seff/(1+4\seff)}$, and the bias is
$2^{-2J_n\seff} \asymp (\log n/n)^{4\seff/(1+4\seff)}$; $p/n$ is of smaller
order.

\emph{Part (iii)(b).} With $s < 1/2$, $\norm{\thetastar}_{1} \asymp
2^{J_n(1/2-s)}$ as in (ii). With
$2^{J_n} \asymp (n/\log n)^{1/(1-s+2\seff)} \gtrsim n/\log n$, because
$\seff \le s/2$, \eqref{eq:lam-plus} gives
$\penn^{+}\norm{\thetastar}_{1} \asymp 2^{J_n(1-s)}\log n/n \asymp
(\log n/n)^{1-(1-s)/(1-s+2\seff)} = (\log n/n)^{2\seff/(1-s+2\seff)}$, which
is also the order of the bias $2^{-2J_n\seff}$. At the boundaries,
$\seff = 1/4$ in (a) and $\seff = s/2$ in (b), $2^{J_n} \asymp n/\log n$,
$\mu_{J_n}$ is bounded, and the rates are $(\log n/n)^{1/2}$ and
$(\log n/n)^{s}$, those of (i) and (ii) at these points.

\emph{Part (iv).} With $s = 1/2$, Lemma S5.4 gives
$\norm{\thetastar}_{1} \le pqC_gJ_n$, and $J_n \asymp \log n$ in both
choices. If $\seff \ge 1/4$,
$2^{J_n} \asymp (n/\log n)^{1/(4\seff)} \lesssim n/\log n$, $\mu_{J_n}$ is
bounded, $\penn^{+}\norm{\thetastar}_{1} \asymp (\log n/n)^{1/2}\log n$,
and the bias, $2^{-2J_n\seff} \asymp (\log n/n)^{1/2}$, is of smaller order.
If $\seff < 1/4$, $2^{J_n} \asymp \{n/(\log n)^{2}\}^{2/(1+4\seff)}$, which
exceeds $n/\log n$ for $n$ large, and $\penn^{+}\norm{\thetastar}_{1} \asymp
2^{J_n/2}(\log n)^{2}/n \asymp \{(\log n)^{2}/n\}^{4\seff/(1+4\seff)}$, the
order of the bias. \qed

% derivations/05-taxas.tex, Observacao rem:teto (E1.14; E5e).
\begin{remark}[Why the resolution passes the ceiling]
\label{rem:ceiling}
In the pairs of (iii) the bound $\penn\norm{\thetastar}_{1} +
\mathcal{B}_{J}^{2}$ is dominated by the bias at every $J$ below the ceiling,
and the best that can be done there is to stop at it: $2^{J_n} = n/(K\log n)$
with $K > 4\norm{\psi}_{\infty}^{2}/(3C_U)$ fixed keeps $\penn$ admissible,
because the exponent of Lemma~\ref{lem:columns} at $t = B_X^{2}C_U$ is then
$-\{3C_UK/(4\norm{\psi}_{\infty}^{2})\}\log n$ and $d_n \le pqn$, and gives
$O_p\{(\log n/n)^{2\seff}\}$. Above the ceiling $\penn^{+}$ pays for the
growth of $\smax$, and the estimation term,
$2^{J/2}(\log n/n)\norm{\thetastar}_{1}$, grows more slowly than the bias
decreases until $J$ reaches the $J_n$ of (iii). The exponents of (iii),
$4\seff/(1+4\seff)$ and $2\seff/(1-s+2\seff)$, are strictly larger than
$2\seff$ away from $\seff = 1/4$ and $\seff = s/2$, and the $J_n$ of (iii)
minimizes the bound over $J$ up to a constant, so that this is the best bound
that the argument gives. In (iv) with $\seff \ge 1/4$, the factor
$(\log n)^{3/2}$ comes from $\norm{\thetastar}_{1} \asymp J_n \asymp \log n$,
which no choice avoids.
\end{remark}

% E5f (E1.15; D54): a prova da Proposicao 8 de derivations/05-taxas.tex.
\subsection*{Proof of Proposition~\ref{prop:lasso-truncated}}

\emph{Part (i).} Lemma~\ref{lem:comparator} with $\bar\bbeta = (\bc,\thbar)$
gives, discarding $\pen\norm{\thetahat}_{1} \ge 0$,
$\normn{\Btil\bar v}^{2} \le 6\pen\norm{\thbar}_{1} + 4\normn{\bbar}^{2}$,
and Step~6 in its proof gives $\normn{\widehat{f}-f} \le \normn{\Btil\bar v}
+ 2\normn{\bbar} + \normn{\PA\bveps}$; with
$(x+y+z)^{2} \le 3(x^{2}+y^{2}+z^{2})$, $\normn{\widehat{f}-f}^{2} \le
18\pen\norm{\thbar}_{1} + 24\normn{\bbar}^{2} + 3\normn{\PA\bveps}^{2}$.
Since $\bbar = \bb + \Bmat(\thetastar-\thbar)$, $\normn{\bbar}^{2} \le
2\normn{\bb}^{2} + 2\normn{\Bmat(\thetastar-\thbar)}^{2}$. If $t$ is
deterministic, so is $u = \thetastar-\thbar$, and
$\E\normn{\Bmat u}^{2} \le \lmaxx(\Gram)\norm{u}_{2}^{2} \le
\kappa_2C_U\norm{u}_{2}^{2}$, as in the proof of
Proposition~\ref{prop:truncated}. With $t = \pen$, on the coordinates with
$|\theta^{*}_{a}| > \pen$ the minimum in \eqref{eq:R1} is
$\pen|\theta^{*}_{a}|$, and on the others it is $\theta^{*2}_{a}$; summing
gives the bound by $\max(18,48\kappa_2C_U)\Rone(\thetastar;\pen)$. The form
in $O_p$ is that of the proof of Proposition~\ref{prop:lasso-slow}: the
event $\mathcal{T}$ with $\penn \ge 2\pen_0$, by Lemma~\ref{lem:columns}
below the ceiling, or with $\penn^{+} \ge 2\pen_0$, by
\eqref{eq:smax-ceiling} for every $J$, and Markov's inequality for
$\normn{\bb}^{2}$, $\normn{\Bmat(\thetastar-\thbar)}^{2}$ and
$\normn{\PA\bveps}^{2}$.

\emph{Part (ii), $s < 1/2$.} Fix a block and a level $\lev$, write
$\theta_{\lev} \in \R^{2^{\lev}}$ for its coefficients, with
$\norm{\theta_{\lev}}_{\pi} \le C_g2^{-\lev(s+1/2-1/\pi)}$ by Assumption~5,
and let $v_{\lev} = \sum_{\tsl}\min(t|\theta_{\lev\tsl}|,\theta_{\lev\tsl}^{2})$.
The first bound is in the $\ell_1$ norm: by H\"older's inequality in
$\R^{2^{\lev}}$, $v_{\lev} \le t\norm{\theta_{\lev}}_{1} \le
t2^{\lev(1-1/\pi)}\norm{\theta_{\lev}}_{\pi} \le tC_g2^{\lev(1/2-s)} =:
A_{\lev}$. The second is the tail, as in the proof of
Proposition~\ref{prop:truncated}(ii) with coordinates in place of chunks:
$v_{\lev} \le t^{2-\pi}\norm{\theta_{\lev}}_{\pi}^{\pi} \le
t^{2-\pi}C_g^{\pi}2^{-\lev\pi\seff}$ if $\pi \le 2$, and
$v_{\lev} \le \norm{\theta_{\lev}}_{2}^{2} \le C_g^{2}2^{-2\lev s}$ if
$\pi \ge 2$; call it $B_{\lev}$. With
\[
  2^{x} = (C_g/t)^{\tau},
  \qquad
  h = C_g^{\tau}t^{2-\tau} = C_g^{\tau}t^{4s/(2s+1)},
\]
one has $A_{\lev} = h2^{(\lev-x)(1/2-s)}$, because $1 + \tau(1/2-s) = \tau$,
and $B_{\lev} = h2^{-(\lev-x)\min(\pi,2)\seff}$, because
$\tau\pi\seff = \pi-\tau$ and $2s\tau = 2-\tau$. Summing $A_{\lev}$ over
$\lev \le x$ and $B_{\lev}$ over $\lev > x$, two geometric series of ratios
$2^{-(1/2-s)}$ and $2^{-\min(\pi,2)\seff}$ that start from at most $h$, gives
at most $h\{(1-2^{s-1/2})^{-1} + (1-2^{-\min(\pi,2)\seff})^{-1}\}$ per block;
truncating at $\lev < J$ only removes terms, and summing over the $pq$
blocks gives \eqref{eq:count-lasso-truncated}. This is the computation of
the proof of Lemma~3, with $\min(ty,y^{2})$ in place of $\min(y^{2},\eta)$.

\emph{Part (ii), $s \ge 1/2$.} The first statement is
$\min(t|y|,y^{2}) \le t|y|$. For the second, the block with
$\theta^{*}_{\cv\md,00} = C_g$ and all its other coefficients zero satisfies
Assumption~5, and gives $\Rone(\thetastar;t) \ge \min(tC_g,C_g^{2}) = tC_g$.
If $s = 1/2$, set, at each level $\lev < J$ with $2^{\lev} \le C_g/t$, all
the $2^{\lev}$ coefficients equal to $C_g2^{-\lev} \ge t$; then
$\norm{\theta_{\lev}}_{\pi} = C_g2^{-\lev(1-1/\pi)}$, the bound of
Assumption~5 at $s = 1/2$, and each such level contributes
$2^{\lev}tC_g2^{-\lev} = tC_g$.

\emph{Part (iii).} With $J_n = \lceil c\log_2 n\rceil$ and $c < 1$,
$2^{J_n}\log d_n/n \asymp n^{c-1}\log n \to 0$, and (i) holds with
$\penn \asymp (\log n/n)^{1/2}$. By \eqref{eq:count-lasso-truncated},
$\Rone(\thetastar;\penn) \lesssim (\log n/n)^{2s/(2s+1)}$. The lower end of
the window \eqref{eq:lasso-window} is $s/\{(2s+1)\seff\}$, because
$(2-\tau)/(4\seff) = s/\{(2s+1)\seff\}$ and $\seff \le s$, so that the bias
is $\mathcal{B}_{J_n}^{2} \lesssim 2^{-2J_n\seff} \le n^{-2c\seff} \le
n^{-2s/(2s+1)}$; and $p/n$ is of smaller order.

\emph{Part (iv).} Write $\eta = 2s/\{s+\seff(2s+1)\}$, which lies in $[1,2)$
exactly when $\seff \le s/(2s+1)$; then $2^{J_n} \asymp (n/\log n)^{\eta}
\gtrsim n/\log n$, $2^{J_n} \le n^{2}$ and, by \eqref{eq:lam-plus},
$\penn^{+} \asymp 2^{J_n/2}\log n/n \asymp (\log n/n)^{1-\eta/2}$. By
\eqref{eq:count-lasso-truncated}, $\Rone(\thetastar;\penn^{+}) \lesssim
(\log n/n)^{(1-\eta/2)4s/(2s+1)}$, and $(1-\eta/2)4s/(2s+1) =
4s\seff/\{s+\seff(2s+1)\}$; the bias is $2^{-2J_n\seff} \asymp
(\log n/n)^{2\seff\eta}$, with the same exponent. At $\seff = s/(2s+1)$,
$\eta = 1$ and the exponent is $2s/(2s+1)$. Outside the window of (iii),
$\pi < 2$, because $\pi \ge 2$ gives $\seff = s > s/(2s+1)$. \qed

\subsection*{Proof of Theorem~\ref{thm:lasso-linear}}

By the definition of $J_n$,
$(n/\log n)^{1/(2\seff+1)} \le 2^{J_n} < 2(n/\log n)^{1/(2\seff+1)}$ and
$J_n \asymp \log n$, so
\[
  \frac{pq2^{J_n}\log(pq2^{J_n})}{n}
  \;\asymp\; \frac{(n/\log n)^{1/(2\seff+1)}\log n}{n}
  \;=\; \Bigl(\frac{\log n}{n}\Bigr)^{2\seff/(2\seff+1)}
  \;\longrightarrow\; 0,
\]
which is the requirement of Proposition 2(ii); \eqref{eq:schur} then gives
$\gtil \ge \gamma/2$ with probability tending to one. On that event
Corollary~\ref{cor:instant} with $s_0 \le d_n$ gives
\[
  \normn{\widehat{f}-f}^2 \;\lesssim\;
  \frac{\sigma^2B_X^2C_U}{\kappa_1c_U}\frac{d_nL_n}{n}
  + \frac{\mathcal{B}^2_{J_n}}{\alpha'} + \frac{\sigma^2p}{n\alpha''} .
\]
With $d_n \asymp pq2^{J_n}$ and $L_n \asymp \log n$ the first term is the
display above; the second is
$\mathcal{B}^2_{J_n} \asymp 2^{-2J_n\seff} \le
(\log n/n)^{2\seff/(2\seff+1)}$ by the lower bound on $2^{J_n}$; and the third
is $O(1/n)$, of strictly smaller order because $2\seff/(2\seff+1)<1$. \qed

\subsection*{Proof of Corollary~\ref{cor:lasso-components}}

Inside a block the wavelets of level below $J$ are orthonormal in $L_2[0,1]$,
so
$\lVert\widehat{g}_{\cv\md}-\PiJ\gcomp{\cv}{\md}\rVert^2_{L_2[0,1]} =
\sum_{\lev\tsl}v^2_{\cv\md,\lev\tsl}$, and summing over the $pq$ blocks gives
$\norm{v}_2^2$; adding the bias of Lemma 2 in each block,
\[
  \sum_{\cv,\md}\lVert\widehat{g}_{\cv\md}-\gcomp{\cv}{\md}
  \rVert^2_{L_2[0,1]} \;\le\; 2\norm{v}_2^2
  + \frac{2pqC_g^2}{1-2^{-2\seff}}2^{-2J\seff} .
\]
Theorem~\ref{thm:lasso-oracle}(ii) bounds $\norm{v}_2^2$ by
$64\pen^2s_0/\gtil^2+16\normn{\bb}^2/\gtil$, and on the event of Theorem~\ref{thm:lasso-linear} the
quantity $\gtil$ is bounded below by the constant $\gamma/2$, so the two terms
are of the same order as those of Theorem~\ref{thm:lasso-linear}, the only difference being a
constant with $\gamma^{-2}$ in place of $\gamma^{-1}$. This is where the
design condition is used a second time. For the levels, Lemma S5.1 gives
$\Amat(\chat-\bc) = \PA(\bb+\bveps)-\PA\Bmat v$, so that
$\normn{\Amat(\chat-\bc)} \le \normn{\bb}+\normn{\PA\bveps}+
\lmaxx(\Gramhat)^{1/2}\norm{v}_2$, while $\Amat\T\Amat/n$ is a principal
submatrix of $\Gramhat$, whence
$\norm{\chat-\bc}_2^2 \le \normn{\Amat(\chat-\bc)}^2/(\gamma/2)$. Finally the
components of a functional coefficient depend on distinct coordinates of $u$
and the basis is orthonormal, so
$\lVert\widehat{\fcoef{\cv}}-\fcoef{\cv}\rVert_{L_2([0,1]^q)} \le
|\widehat{c}_{\cv}-\lvl{\cv}| + \sum_{\md}\lVert\widehat{g}_{\cv\md}-
\gcomp{\cv}{\md}\rVert_{L_2[0,1]}$, with $q$ fixed. The same argument applies
verbatim under the conditions of Theorem~\ref{thm:lasso-main}. \qed

\subsection*{Proof of Theorem~\ref{thm:lasso-main}}

\emph{A nonasymptotic bound.} Apply Lemma~\ref{lem:comparator} with
$\bar\bbeta = (\bc,\thstarM{M})$, the best $M$ term approximation of the
oracle. Then $\bar M \le M$ and
$\bbar = \bb + \Bmat(\thetastar-\thstarM{M})$, so that, by the
Cauchy-Schwarz inequality, Lemma~\ref{lem:lasso-weak} and \eqref{eq:lmax},
\[
  \normn{\bbar}^2 \;\le\; 2\normn{\bb}^2
  + 2\lmaxx(\Gramhat)\norm{\thetastar-\thstarM{M}}_2^2
  \;\le\; 2\normn{\bb}^2 + \frac{2\Lambda C_\tau^2\tau}{2-\tau}M^{1-2/\tau} .
\]
Substituting in \eqref{eq:oracle-min}, on the intersection of $\mathcal{T}$
with $\{\lmin(\Gramhat)\ge\gamma/2\}$,
\[
  \normn{\widehat{f}-f}^2 \;\le\; 3\min_{1\le M\le d}h(M)
  + 120\normn{\bb}^2 + 3\normn{\PA\bveps}^2,
  \qquad
  h(M) = \frac{64\pen^2}{\gtil}M
  + \frac{40\Lambda C_\tau^2\tau}{2-\tau}M^{1-2/\tau} .
\]
The function $h$ is strictly convex on the positive half line, since
$1-2/\tau<0$, and $h'$ vanishes at
$\Mstar = \{5\Lambda C_\tau^2\gtil/(8\pen^2)\}^{\tau/2}$, where
$h(\Mstar) = 2(2-\tau)^{-1}(64\pen^2/\gtil)\Mstar$; moreover
$h(\lceil\Mstar\rceil) \le h(\Mstar)+64\pen^2/\gtil \le 2h(\Mstar)$ whenever
$\Mstar \ge 1$, because $(2-\tau)/2<1$.

\emph{The resolution.} With $\penn^2 \asymp \sigma^2B_X^2C_UJ_n/n$,
$\gtil \in [\gamma/2,\Lambda]$ and the other constants fixed,
$\Mstar_n \asymp \penn^{-\tau} \asymp (n/J_n)^{\tau/2}$ and
$h(\Mstar_n) \asymp \penn^2\Mstar_n \asymp (J_n/n)^{(2-\tau)/2}$. Writing
$J_n = \lceil c\log_2 n\rceil$, so that $2^{J_n} \asymp n^{c}$, the three
inequalities of the window \eqref{eq:lasso-window} are, in order: $\Mstar_n \le d_n$, which asks
$(n/J_n)^{\tau/2} \lesssim pq\,n^{c}$, that is $c \ge \tau/2$ up to
logarithms; that the bias not dominate, $n^{-2c\seff} \lesssim
(J_n/n)^{(2-\tau)/2}$, that is $c \ge (2-\tau)/(4\seff)$; and the requirement
of Proposition 2(ii), $n^{c-1}\log n \to 0$, that is $c<1$. The term
$\sigma^2p/n$ is of smaller order because $(2-\tau)/2<1$. The window is
nonempty if and only if $\tau/2<1$, which always holds, and
$(2-\tau)/(4\seff)<1$, that is $\seff>(2-\tau)/4$.

\emph{Under Assumption 5 the hypothesis is automatic.} By Lemma~\ref{lem:lasso-weak} the
envelope holds with $\tau = (s+1/2)^{-1}$, whence $2-\tau = 4s/(2s+1)$,
$(2-\tau)/2 = 2s/(2s+1)$ and $(2-\tau)/4 = s/(2s+1)$, which is the condition
$\seff>s/(2s+1)$ of the statement; for $\pi \ge 2$ it holds for every $s>0$,
and for $\pi<2$ it reads $2s^2/(2s+1)>1/\pi-1/2$. The rate for the components
is the argument of Corollary~\ref{cor:lasso-components} with
$\norm{\thetahat-\thetastar}_2 \le \norm{\bar v}_2 +
\norm{\thetastar-\thstarM{M}}_2$ and Lemma~\ref{lem:comparator}, the only difference being an
extra factor $\gtil^{-1}$ in the first term, which does not change the order.
Finally $2s/(2s+1) \ge 2\seff/(2\seff+1)$, with equality only when
$\seff = s$, by monotonicity of $x \mapsto 2x/(2x+1)$ and $\seff \le s$, so
the rate is never worse than that of Theorem~\ref{thm:lasso-linear} and is strictly better when
$\pi<2$. \qed

\begin{remark}[What compressibility buys]
\label{rem:s0}
Read at the oracle, Theorem~\ref{thm:lasso-oracle}(ii) has $s_0 = d$ and delivers the rate of
linear approximation; read at a sparse comparator, through Lemma~\ref{lem:comparator}, it delivers the rate
of Theorem~\ref{thm:lasso-main}. It is the same inequality in both cases, and neither says
anything about $\supp(\thetahat)$: no minimal signal and no
irrepresentability condition is used, and what is gained is a rate governed by
the effective dimension $\Mstar_n$ rather than by $d_n$.
\end{remark}

\subsection*{Proof of Corollary~\ref{cor:lasso-threshold}}

\emph{A nonasymptotic bound on $\Delta$.} Work at $J = J_n$ and $\pen = \penn$,
and intersect the four events of the proof of Corollary~\ref{cor:instant}: $\mathcal{T}$,
of probability at least $1-\alpha$, by Lemma~\ref{lem:calib};
$\{\smax^2 \le 2B_X^2C_U\}$, by Lemma S5.3, which makes $\penn$ admissible;
$\{\lmin(\Gramhat) \ge \gamma/2\}$, by Proposition 2(ii), which with
\eqref{eq:schur} gives $\gtil \ge \gamma/2$; and
$\{\normn{\bb}^2 \le \mathcal{B}^2_{J_n}/\alpha'\}$, by Markov's inequality and
Lemma 2. On this event Theorem~\ref{thm:lasso-oracle}(ii) with $s_0 \le d_n$ gives
\[
  \norm{v}_2^2 \;\le\; \frac{64\penn^2d_n}{\gtil^2}
  + \frac{16\normn{\bb}^2}{\gtil}
  \;\le\; \frac{256\penn^2d_n}{\gamma^2}
  + \frac{32\mathcal{B}^2_{J_n}}{\alpha'\gamma} .
\]
For a fixed block, the orthonormality of the wavelets of level below $J_n$
gives $\lVert\widehat{g}_{\cv\md} - \PiJ\gcomp{\cv}{\md}\rVert_{L_2[0,1]} =
\norm{v_{\cv\md,\cdot}}_2 \le \norm{v}_2$, and the first display of Lemma 2
bounds $\lVert\PiJ\gcomp{\cv}{\md} - \gcomp{\cv}{\md}\rVert^2_{L_2[0,1]}$ by
$C_g^2(1-2^{-2\seff})^{-1}2^{-2J_n\seff}$. By $(x+y)^2 \le 2x^2 + 2y^2$,
\begin{equation}
  \label{eq:Dn}
  \Delta^2 \;\le\; \bar{\Delta}_n^2 \;=\; \frac{512\,\penn^2d_n}{\gamma^2}
  + \frac{64\,\mathcal{B}^2_{J_n}}{\alpha'\gamma}
  + \frac{2C_g^2}{1-2^{-2\seff}}\,2^{-2J_n\seff}
\end{equation}
with probability at least
$1 - \alpha - \alpha' - \Prob\{\lmin(\Gramhat) < \gamma/2\} -
\Prob\{\smax^2 > 2B_X^2C_U\}$, the last two terms tending to zero by
Proposition 2(ii) and Lemma S5.3; their sum is the
$\omega_n$ of Corollary~\ref{cor:lasso-threshold}(iii), and this paragraph is the proof of that
part. By Lemma S7.1(iii), on that event
$\widehat{\mathcal{S}}(t) = \mathcal{S}$ for every $t$ with $\bar{\Delta}_n \le t < \nu_{\min} - \bar{\Delta}_n$, a window
that is nonempty whenever $\nu_{\min} > 2\bar{\Delta}_n$. The bound $\bar{\Delta}_n$ depends on the
components only through the constants of Assumption 5, and the two residual
probabilities do not depend on them at all, so everything above holds
uniformly over the class, which is what allows the target to vary with $n$ in
Assumption 6.

\emph{The order of $\bar{\Delta}_n$.} At the resolution of Theorem~\ref{thm:lasso-linear} the three terms of
\eqref{eq:Dn} are at most of order $\rho_n^2$: the first because
$\penn^2d_n \asymp pq2^{J_n}J_n/n$, which is the quantity shown to be of order
$(\log n/n)^{2\seff/(2\seff+1)}$ in the proof of Theorem~\ref{thm:lasso-linear}, and the other two
because $\mathcal{B}^2_{J_n} \asymp 2^{-2J_n\seff} \le
(\log n/n)^{2\seff/(2\seff+1)}$ by the definition of $J_n$. Hence
$\bar{\Delta}_n \le c(\alpha,\alpha')\rho_n$ with $c(\alpha,\alpha') < \infty$. Given $\epsilon > 0$,
take $\alpha = \alpha' = \epsilon/4$ and $n$ large enough that the two
residual probabilities add up to less than $\epsilon/2$; then
$\Prob\{\Delta > c(\epsilon/4,\epsilon/4)\rho_n\} < \epsilon$, that is
$\Delta = O_p(\rho_n)$.

\emph{Part (i).} If $t_n/\rho_n \to \infty$, then
$\Prob(\Delta \le t_n) \to 1$, and on that event Lemma S7.1(i) with
$\bar{\Delta} = \Delta$ gives $\widehat{\mathcal{S}}(t_n) \subseteq \mathcal{S}$. No property of $\nu_{\min,n}$
has been used.

\emph{Part (ii).} If $\mathcal{S}$ is empty, part (ii) is part (i). Otherwise, the
condition $\limsup_n t_n/\nu_{\min,n} < 1$ gives $\eta > 0$ and $n_0$ with
$\nu_{\min,n} - t_n \ge \eta\,\nu_{\min,n}$ for $n \ge n_0$, and Assumption 6 with $\rho_n$ in place of $\rhoG$ gives
$(\nu_{\min,n} - t_n)/\rho_n \ge \eta\,\nu_{\min,n}/\rho_n \to \infty$. Hence
$\Prob(\Delta \le t_n) \to 1$ and $\Prob(\Delta < \nu_{\min,n} - t_n) \to 1$, and
on the intersection of the two events Lemma S7.1(iii) with $\bar{\Delta} = \Delta$ gives
$\widehat{\mathcal{S}}(t_n) = \mathcal{S}$. For the existence, $t_n = (\rho_n\nu_{\min,n})^{1/2}$ has
$t_n/\rho_n = (\nu_{\min,n}/\rho_n)^{1/2} \to \infty$ and
$t_n/\nu_{\min,n} = (\rho_n/\nu_{\min,n})^{1/2} \to 0$. \qed

\emph{At the resolution of Theorem~\ref{thm:lasso-main}.} Theorem~\ref{thm:lasso-main} gives
$\sum_{\cv,\md}\lVert\widehat{g}_{\cv\md}-\gcomp{\cv}{\md}\rVert^2_{L_2[0,1]}
= O_p(\tilde\rho_n^{\,2})$ with $\tilde\rho_n = (\log n/n)^{s/(2s+1)}$, hence
$\Delta = O_p(\tilde\rho_n)$, and the bound in its proof depends on the
components only through the constants of Assumption 5, through $C_\tau$ of
Lemma~\ref{lem:lasso-weak} and the design constants, so that this holds uniformly over the
class. The proofs of parts (i) and (ii) above use nothing about $\Delta$
beyond $\Delta = O_p(\rho_n)$ uniformly over the class, together with
Lemma S7.1 and Assumption 6, and they go through verbatim with
$\tilde\rho_n$ in place of $\rho_n$.


%%%%%%%%%%%%%%%%%%%%%%%%%%%%%%%%%%%%%%%%%%%%%%%%%%%%%%%%%%%%%%%%%%%%%%%%%%%%%%
% E5h (k = 5): a S9, o que completa a Secao 5 do ms_5 (wafc/R/dgp.R,
% wafc_component() e wafc_window(); wafc-studies/config/scale.yaml), e os lugares
% das tabelas do estudo, sem numeros.

